%\setcounter{definition}{0} \setcounter{property}{0} \setcounter{claim}{0} \setcounter{fact}{0} \setcounter{corollary}{0} \setcounter{figure}{0}
\section{Edit Distance Problem}

The edit distance problem is largely motivated by and abstracted from genome evolution.
There are three types of nucleotide-level mutations, substitutions, insertions, and deletions.
See Figure~\ref{fig:edits}. Given two DNA sequences, for examples, the sequences of two homologous genes,
we want to infer the most likely mutation events happened between them during the evolutionary history.
This can be modeled as the following edit distance problem: given two strings $X$ and $Y$ over alphabet $\Sigma$,
to find the minimum number of such events that transforms $X$ into $Y$.
We denote by $d(X, Y)$ as the edit distance between $X$ and $Y$.
See Figure~\ref{fig:align} for an example, where we can write $d(X, Y) = 4$.

\begin{figure}[h]
\centering{\input{edits}}
\caption{Illustration of the 3 types of mutation events in genome evolution.}
\label{fig:edits}
\end{figure}

An equivalent and convenient way to represent edits between two sequences 
is \emph{alignment}. See Figure~\ref{fig:align}.
In an alignment, deletions and insertions are represented with \emph{gaps}, shown with $-$.
(Note that deletions on $X$ is equivalent to insertions on $Y$;
deletions on $Y$ is equivalent to insertions on $X$.)
Both exactly-matched characters and mismatched-characters~(i.e., requires a substitutions event),
are shown by vertical-bars with substitutions highlighted.
We define the \emph{score} of an alignment as the number of mismatches
and the gaps in $X$ and the gaps in $Y$.
Clearly, there is a one-to-one correspondence between a series
of events that transforms $X$ into $Y$ and an alignment between $X$ and $Y$,
and the number of events is equal to the score of the alignment. %the number of gaps and mismatches in the alignment.
Hence, calculating edit distance between $X$ and $Y$ is equivalent to
seeking the optimal alignment between $X$ and $Y$~(i.e., the alignment with minimized score).

\begin{figure}[h]
\centering{\input{align}}
\caption{Illustration of the edit distance and the corresponding optimal alignment.}
\label{fig:align}
\end{figure}


\subsection*{A Dynamic Programming Algorithm}

We now design a dynamic programming to find the optimal alignment of two given strings.
We first look into the optimal substructures. Suppose that $\mathcal{A}$ is the optimal
alignment of two strings $X$ and $Y$~(see Figure~\ref{fig:optimal}).
Then clearly its substructure, i.e., the first $k$ columns, gives the optimal
alignment between the corresponding two substrings, which can be obtained by removing
the gaps in the first $k$ columns of $\mathcal{A}$.
(Proof: if the two substrings can be aligned with
a better alignment~(with smaller score), then we can have
a better alignment, than $\mathcal{A}$, between $X$ and $Y$,
by concatenating to it the remaining columns of $\mathcal{A}$.)

\begin{figure}[h]
\centering{\input{optimal}}
\caption{Illustration of the optimal substructure of alignment.
The left figure shows the optimal alignment between $X$ and $Y$.
Then we can remove any number of columns from right-side, and the resulting
alignment must be the optimal alignment of the corresponding two substrings.
For examples, the middle figure shows the optimal alignment between $X' = AGCATCAG$ and $Y' = AGGAAG$,
where one column of the alignment in the left-figure is removed;
and the right figure shows the optimal alignment between $X'' = AGCATC$ and $Y'' = AGGA$,
where three columns of the alignment in the left-figure is removed.}
\label{fig:optimal}
\end{figure}

So this problem satisfies the optimal substructure property.
The substructure suggests defining the subproblems as to computing the optimal
alignment between a substring of $X$~(starting from beginning) and 
a substring of $Y$~(starting from beginning).
Formally, let $F(i,j)$ be the edit distance between $X[1\cdots i]$ and $Y[1\cdots j]$.
Clearly, $d(X, Y) = F(m,n)$, where we assume $m = m$ and $n = n$.
Now we develop a recurrence. In order to calculate the optimal alignment between $X[1\cdots i]$
and $Y[1\cdots j]$, consider the possible scenario of the last step, i.e., how $X[i]$ and $X[j]$
are placed in the optimal alignment. There are three cases. See Figure~\ref{fig:recurrence}.

\begin{figure}[h]
\centering{\input{recurrence}}
\caption{Illustration of the three cases in developing the recurrence.}
\label{fig:recurrence}
\end{figure}


First, $X[i]$ is aligned to $Y[j]$ in the optimal alignment. In this case, the optimal alignment
between $X[1\cdots i]$ and $Y[1\cdots j]$ can be obtained by concatenating the optimal alignment
between $X[1\cdots i-1]$ and $Y[1\cdots j-1]$ and this column of $(X[i], Y[j])^T$.
Hence, in this case, $F(i,j) = F(i-1, j-1) + \delta(X[i] \neq Y[j])$, where the delta function
$\delta(X[i] \neq Y[j]) = 1$ if $X[i] \neq Y[j]$, i.e., a mismatch that requires a substitution event,
and $\delta(X[i] \neq Y[j]) = 0$ if $X[i] = Y[j]$, i.e., an exact-match that does not need any event.
%The extra delta function represents if a substitution is needed.

Second, $X[i]$ is not aligned~(or, aligned to a gap) in the optimal alignment. In this case, the optimal alignment
between $X[1\cdots i]$ and $Y[1\cdots j]$ can be obtained by concatenating the optimal alignment
between $X[1\cdots i-1]$ and $Y[1\cdots j]$ and this column of $(X[i], -)^T$.
Hence, in this case, $F(i,j) = F(i-1, j) + 1$, where 
the extra 1 represents a deletion on $X$ is needed.

Third, $Y[j]$ is not aligned~(or, aligned to a gap) in the optimal alignment. In this case, the optimal alignment
between $X[1\cdots i]$ and $Y[1\cdots j]$ can be obtained by concatenating the optimal alignment
between $X[1\cdots i]$ and $Y[1\cdots j-1]$ and this column of $(-, Y[j])^T$.
Hence, in this case, $F(i,j) = F(i, j- 1) + 1$, where 
the extra 1 represents a deletion on $Y$ is needed.

Combined, the recurrence is given below.
\begin{displaymath}
F(i,j) = \min\left\{
	\begin{array}{llll}
	F(i-1,j-1) + \delta(X[i] \neq Y[j]) \\
	F(i-1,j) + 1 \\
	F(i,j-1) + 1 \\
	\end{array}
\right.
\end{displaymath}

We now can complete the algorithm. The algorithm fills a DP table. See Figure~\ref{fig:table}.
To facilitate the use of the recurrence to solve all subproblems, we add a new row and a new column in the table,
and fill them in the initialization step. 
We then fill the table, either row by row, or column by column. (Think: why both work?).
The right bottom entry, $F(m, n)$, gives the edit distance between $X$ and $Y$.
The actual alignment can be constructed by maintaining a tracking back pointer on each entry.

\begin{figure}[t!]
\centering{\input{table}}
\caption{The dynamic programming table.}
\label{fig:table}
\end{figure}

The runtime of this algorithm is $\Theta(mn)$.
There are faster algorithm to calculate the edit distance~(and the optimal alignment).
For example, using the Four Russian's technique, it can be solved in $O(n^2/\log n)$, where we assume that $m = n$.
However, it is unlikely to design strong subquadratic algorithms.
Formally, it is proved that, unless the exponential time hypothesis is false, there does not exist $O(n^{2-\epsilon})$ time
algorithm for any $\epsilon > 0$.


\begin{minipage}{0.8\textwidth}
	\aaA {13}{Algorithm DP-optimal-alignment~(two strings $X$, $Y$)}\xxx
	\aab {$F(0,j) = j$, for any $0 \le j \le n$}\xxx
	\aab {$F(i,0) = i$, for any $0 \le i \le m$}\xxx
	\aab {$T(0,j) = null$, for any $0 \le j \le n$}\xxx
	\aab {$T(i,0) = null$, for any $0 \le i \le m$}\xxx
	\aaB {5}{for $i = 1 \to m$}\xxx
	\aaC {3}{for $j = 1 \to n$}\xxx
	\aad {$F(i,j) = \min\{F(i-1,j-1) + \delta(X[i] \neq Y[j]), F(i-1,j) + 1, F(i,j-1) + 1\}$;}\xxx
	\aad {let $T(i,j)$ store the best case}\xxx
	\aac {end for;}\xxx
	\aab {end for;}\xxx
	\aab {report: $F(m,n) = d(X,Y)$;}\xxx
	\aab {the optimal alignment can be obtained by tracing back following $T$;}\xxx
	\aaa {end algorithm;}\xxx
\end{minipage}


\subsection*{Transforming into Shortest Path Problem}

Another way to solve the edit distance problem is to transform it into a shortest-path problem.
See Figure~\ref{fig:graph}.
We build a graph with $(|X| + 1)(|Y| + 1)$ vertices, $v_{ij}$, $0\le 0 \le |X|$ and $0\le j \le |Y|$. 
(Imagine placing a vertex in each entry of above DP table.)
There are three in-edges for vertex $v_{ij}$, when $i\ge 1$ and $j \ge 1$, corresponding to
the 3 cases of the recursion. 
All vertical edges have edge length of 1;
all horizontal edges have edge length of 1;
all diagonal edges have edge length of either 1 or 0: 
if $X[i] \neq Y[j]$ then the edge $(v_{i-1,j-1}, v_{ij})$ has a length of 1,
and otherwise this edge has a length of 0.

\begin{figure}[h]
\centering{\input{graph}}
\caption{The graph constructed for $X=ACG$ and $Y = AGGT$. The 3 blue edges have length of 0, and all other edges have length 1.}
\label{fig:graph}
\end{figure}

Clearly, there is a one-to-one correspondence between a path from $v_{00}$ to $v_{|X|,|Y|}$
and an alignment between $X$ and $Y$, and the length of the path equals to the
number of operations in the alignment. 
See Figure~\ref{fig:graph2}.
In fact, a path from $v_{00}$ to $v_{|X|,|Y|}$
suggests how to align $X$ and $Y$: a diagonal edge pointing to vertex $v_{ij}$ implies aligning $X[i]$ and $Y[j]$;
a vertical edge pointing to vertex $v_{ij}$ implies not to aligning $X[i]$~(or aligning $X[i]$ with gap);
a horizontal edge pointing to vertex $v_{ij}$ implies not to aligning $Y[j]$~(or aligning $Y[j]$ with gap).
You can also see that the length of this path equals to the number of 
gaps and mismatches in the resulting alignment. This correlation is bi-directional:
an alignment can be mapped into a path in the graph too.


\begin{figure}[h!]
\centering{\input{graph2}}
\caption{Two examples showing the one-to-one correspondence between paths in the graph and alignments.}
\label{fig:graph2}
\end{figure}

The above analysis leads to that, the optimal alignment between $X$ and $Y$
corresponds to the shortest path from  $v_{00}$ to $v_{|X|,|Y|}$.
As the graph is a DAG, we can use the linear time algorithm to find the shortest path from $v_{00}$.
The graph contains $\Theta(|X||Y|)$ edges and vertices.
Hence the running time is $\Theta(|X||Y|)$,
the same as the DP algorithm.


\subsection*{Why not Recursive Programming?}

In the two DP algorithms we've designed, we always use a table to store
the optimal solution of subproblems; when solving larger
subproblems, we simply query the table and fetch the optimal
solutions as needed. 
This leads to a bottom-up strategy, i.e., we start from solving
smaller subproblems as they will be needed by larger ones.
Is this strategy always necessary?
Can we directly solve the original problem with recursive programming?

Let's try.
Consider the edit distance problem.
The DP algorithm we designed there uses the recursion below.
\begin{displaymath}
F(i,j) = \min\left\{
	\begin{array}{llll}
	F(i-1,j-1) + \delta(X[i] \neq Y[j]) \\
	F(i-1,j) + 1 \\
	F(i,j-1) + 1 \\
	\end{array}
\right.
\end{displaymath}

We write a recursive procedure: define $F(i,j)$
returns the edit distance between $X[1\cdots i]$ and $Y[1\cdots j]$.

\begin{minipage}{0.8\textwidth}
	\aaA {7}{function $F(i,j)$}\xxx
	\aab {if $i = 0$: return $j$;}\xxx
	\aab {if $j = 0$: return $i$;}\xxx
	\aab {$z_1 = F(i-1,j-1) + \delta(X[i] \neq Y[j])$;}\xxx
	\aab {$z_2 = F(i-1,j) + 1$;}\xxx
	\aab {$z_3 = F(i,j-1) + 1$;}\xxx
	\aab {return $\min\{z_1,z_2,z_3\}$;}\xxx
	\aaa {end function;}\xxx
\end{minipage}

We then call $F(m,n)$, where $m=|X|$ and $n=|Y|$.
This is a correct algorithm.
How about its running time?
We define $T(m,n)$ as the running time of $F(m,n)$.
We have $T(m,n) = T(m-1,n-1) + T(m-1,n) + T(m,n-1) + 1$.
To \emph{approximate} $T(m,n)$, assume that $m=n$.
Notice also that $T(m-1,n) \ge T(m-1,n-1)$
and $T(m,n-1) \ge T(m-1,n-1)$.
We can write $T(n,n) \ge 3T(n-1,n-1) + 1$.
Hence, $T(n,n) \ge \Theta(3^n)$.
So, above algorithm runs in exponential-time!
This suggests that recursive procedure is a bad strategy.

Note that there are $\Theta(mn)$ distinct subproblems, i.e., $F(i,j)$, $0\le i
\le m$, $0\le j \le n$. The reason why above
algorithm runs in exponential time is that it keeps solving the \emph{same}
subproblems. See Figure~\ref{fig:search-tree} for the recursive tree. 
The height of that tree is linear, i.e., $\Omega(n)$~(again assuming $m=n$), as
the size of the subproblem is usually reduced by a \emph{constant}~(for example, $n\to n-1$) in each recursive call.
The total number of nodes in the tree is $\Omega(3^n)$ which is exponential
and aligns with above runtime analysis.
But there are only $\Theta(mn)$ distinct subproblems, showing that
vast duplicates of nodes in the tree.
Hence, in the DP it is necessary to use a table to store the optimal solutions so as to avoid
repeating resolving the same ones.

\begin{figure}[h]
\centering{\input{search-tree}}
\caption{First 3 levels of the recursive tree of above procedure.
Identical recursive calls are marked with the same color.}
\label{fig:search-tree}
\end{figure}




%An alternative way to understand it is that, subproblems will be used
%multiple times by different larger subproblems, hence it's beneficial
%to store them.

Both divide-and-conquer algorithms and dynamic programming algorithm
rely on a recursive formula that reduces larger subproblems into smallers.
See the merge-sort algorithm given below.

\begin{minipage}{0.8\textwidth}
        \aaA {5}{Algorithm merge-sort~($S[1\cdots n]$)}\xxx
        \aab {if $n \le 1$: return $S$;}\xxx
        \aab {$S'_1 =$ merge-sort~($S[1\cdots n/2]$);}\xxx
        \aab {$S'_2 =$ merge-sort~($S[n/2 + 1\cdots n]$);}\xxx
        \aab {return merge-two-sorted-arrays~($S_1', S_2'$);}\xxx
        \aaa {end algorithm;}\xxx
\end{minipage}

Why for D\&C algorithm we do not maintain a table?
There are two reasons.
First, in (most) D\&C algorithms including above merge-sort algorithm, 
there are no duplicate subproblems~(e.g., in merge-sort all subproblems are distinct).
Hence, there is no difference between
solving recursively or by using a table as they solve the same amount of subproblems.
Second, in (most) D\&C algorithms, the size of the subproblem is usually reduced by a \emph{fixed ratio}~(for example, $n\to n/2$),
and consequently the height of the recursive tree is $\Theta(\log n)$ and there is linear number of nodes in the recursive tree. 
Hence, there is \emph{no need} to maintain a table, as solving all subproblems recursively is still efficient.

